\documentclass[11pt, a4paper]{article}

\usepackage[T1]{fontenc}
\usepackage[utf8]{inputenc}
\usepackage[spanish]{babel}
\usepackage[a4paper, margin=2.5cm]{geometry}
\usepackage{amsmath, amssymb}
\usepackage{booktabs}
\usepackage{enumitem}
\usepackage{listings}
\usepackage{xcolor}
\usepackage{hyperref}

\definecolor{codigobg}{RGB}{245,245,245}
\definecolor{keyword}{RGB}{0,100,150}
\lstset{language=R, basicstyle=\ttfamily\small, backgroundcolor=\color{codigobg},
  keywordstyle=\color{keyword}\bfseries, frame=single, breaklines=true, showstringspaces=false}

\setlength{\parindent}{0pt}
\pagestyle{plain}

\begin{document}

\begin{center}
{\large\bfseries PONTIFICIA UNIVERSIDAD JAVERIANA}\\[3pt]
Facultad de Ciencias Econ\'omicas y Administrativas\\[2pt]
Estad\'istica 2026-II --- Prof.\ Jaime Polanco Jim\'enez\\[6pt]
{\LARGE\bfseries Problem Set 8}\\[3pt]
{\large Regresi\'on Lineal M\'ultiple y Diagn\'osticos Avanzados}\\[6pt]
Asignado: Octubre 30 (Sesi\'on 10) \quad$|$\quad Entrega: Noviembre 13 (Sesi\'on 12)\\[3pt]
\textit{Plazo extendido: 2 semanas}
\end{center}

\vspace{6pt}
\rule{\linewidth}{0.5pt}
\vspace{2pt}
{\centering\small Repositorio: \url{https://github.com/polanco-jaime/Curso-Estadistica-2026-3}\par}
\vspace{6pt}

\textbf{Instrucciones.} Este Problem Set es individual y se resuelve en R sobre
microdatos ICFES. La entrega es digital al inicio de la clase indicada.
El c\'odigo R debe incluirse. Este PS integra los conceptos de regresi\'on
m\'ultiple y diagn\'osticos avanzados como cierre del curso.

\vspace{10pt}

\section*{Enunciado}

En este Problem Set (capstone) estimar\'a un modelo de regresi\'on lineal
m\'ultiple para predecir el puntaje de ingl\'es en Saber Pro de estudiantes
de Negocios Internacionales (NI), usando m\'ultiples regresores de Saber 11,
y realizar\'a un diagn\'ostico exhaustivo del modelo.

\begin{enumerate}[label=\arabic*)]
  \item \textbf{Modelo OLS m\'ultiple con regresores de Saber 11:} estime un
        modelo para predecir \texttt{MOD\_INGLES\_PUNT} (Saber Pro) usando
        \textbf{al menos 4 regresores} del examen Saber 11:
        \begin{enumerate}[label=\alph*)]
          \item Seleccione 4 o m\'as variables de Saber 11 como regresores.
                Sugerencias:
                \begin{itemize}
                  \item \texttt{PUNT\_INGLES} (ingl\'es)
                  \item \texttt{PUNT\_MATEMATICAS} (matem\'aticas)
                  \item \texttt{PUNT\_LECTURA\_CRITICA} (lectura cr\'itica)
                  \item \texttt{PUNT\_SOCIALES\_CIUDADANAS} (sociales y
                        ciudadanas)
                  \item \texttt{PUNT\_CIENCIAS\_NATURALES} (ciencias naturales)
                \end{itemize}
          \item Estime el modelo OLS:
                \[
                  \texttt{MOD\_INGLES\_PUNT}_i = \beta_0 + \beta_1 X_{1i} +
                  \beta_2 X_{2i} + \cdots + \beta_k X_{ki} + \varepsilon_i
                \]
                usando \texttt{lm()} y presente el \texttt{summary()} completo.
          \item Interprete cada coeficiente $\hat{\beta}_j$ en contexto.
                ¿Cu\'al regresor tiene el mayor efecto sobre el puntaje de
                ingl\'es en Saber Pro?
        \end{enumerate}

  \item \textbf{Multicolinealidad: Factor de Inflaci\'on de la Varianza (VIF):}
        eval\'ue si existe multicolinealidad entre los regresores.
        \begin{enumerate}[label=\alph*)]
          \item Calcule el VIF para cada regresor usando \texttt{car::vif()}.
          \item Reporte los valores de VIF. ¿Hay alg\'un regresor con
                $\text{VIF} > 5$ o $\text{VIF} > 10$?
          \item Si hay multicolinealidad alta, ¿qu\'e implica para la
                precisi\'on de las estimaciones de los coeficientes?
          \item ¿Considera necesario eliminar alg\'un regresor del modelo?
                Justifique.
        \end{enumerate}

  \item \textbf{Sesgo por variable omitida (OVB):} compare el coeficiente de
        \texttt{PUNT\_INGLES} en dos modelos para ilustrar el sesgo por
        omisi\'on de variables:
        \begin{enumerate}[label=\alph*)]
          \item Estime un modelo simple:
                \[
                  \texttt{MOD\_INGLES\_PUNT}_i = \gamma_0 + \gamma_1
                  \texttt{PUNT\_INGLES}_i + u_i
                \]
          \item Compare $\hat{\gamma}_1$ (modelo simple) con
                $\hat{\beta}_1$ (modelo m\'ultiple del punto 1).
          \item ¿Cambi\'o el coeficiente? ¿En qu\'e direcci\'on (aument\'o o
                disminuy\'o)?
          \item Explique el sesgo por variable omitida: ¿por qu\'e el modelo
                simple sobrestima o subestima el efecto de
                \texttt{PUNT\_INGLES}?
        \end{enumerate}

  \item \textbf{Diagn\'osticos gr\'aficos: 4 gr\'aficos base:} produzca los
        4 gr\'aficos diagn\'osticos est\'andar de \texttt{plot(modelo)} para
        el modelo m\'ultiple:
        \begin{enumerate}[label=\alph*)]
          \item Residuos vs.\ valores ajustados
          \item Q-Q plot de residuos
          \item Scale-Location (ra\'iz de residuos estandarizados vs.\
                ajustados)
          \item Residuos vs.\ leverage (con curvas de distancia de Cook)
        \end{enumerate}
        Interprete cada gr\'afico: ¿qu\'e supuesto eval\'ua (linealidad,
        homocedasticidad, normalidad, observaciones influyentes) y qu\'e
        concluye?

  \item \textbf{Prueba de heterocedasticidad (Breusch-Pagan):} eval\'ue si
        los errores tienen varianza constante.
        \begin{enumerate}[label=\alph*)]
          \item Plantee las hip\'otesis:
                \[
                  H_0: \text{homocedasticidad} \quad \text{vs.} \quad
                  H_1: \text{heterocedasticidad}
                \]
          \item Aplique la prueba de Breusch-Pagan usando
                \texttt{lmtest::bptest()}.
          \item Reporte el estad\'istico de prueba y el $p$-valor.
          \item ¿Rechaza $H_0$ al 5\%? ¿Qu\'e implica para los errores
                est\'andar OLS si hay heterocedasticidad?
        \end{enumerate}

  \item \textbf{Observaciones influyentes: distancia de Cook:} identifique
        observaciones que tienen influencia desproporcionada en el modelo.
        \begin{enumerate}[label=\alph*)]
          \item Calcule la distancia de Cook para cada observaci\'on usando
                \texttt{cooks.distance()}.
          \item Identifique las 5 observaciones m\'as influyentes (mayor
                distancia de Cook).
          \item ¿Alguna observaci\'on supera el umbral $D_i > 4/n$?
          \item Examine esas observaciones: ¿qu\'e caracter\'isticas tienen?
                ¿Son outliers genuinos o datos at\'ipicos v\'alidos?
        \end{enumerate}

  \item \textbf{Errores est\'andar robustos (HC1):} si detect\'o
        heterocedasticidad, reestime los errores est\'andar usando la
        correcci\'on de Huber-White.
        \begin{enumerate}[label=\alph*)]
          \item Use \texttt{sandwich::vcovHC()} con \texttt{type = "HC1"}
                para calcular la matriz de varianza-covarianza robusta.
          \item Reporte los coeficientes con errores est\'andar robustos
                usando \texttt{lmtest::coeftest()}.
          \item Compare los errores est\'andar OLS (del punto 1) con los
                robustos. ¿Cu\'ales son m\'as grandes?
          \item ¿Cambian las conclusiones de significancia estad\'istica para
                alg\'un regresor?
        \end{enumerate}

  \item \textbf{Bondad de ajuste y comparaci\'on de modelos:} eval\'ue el
        poder explicativo del modelo.
        \begin{enumerate}[label=\alph*)]
          \item Reporte $R^2$ y $R^2$ ajustado del modelo m\'ultiple.
          \item Compare el $R^2$ del modelo m\'ultiple con el $R^2$ del modelo
                simple del punto 3. ¿Cu\'anto aumenta el poder explicativo al
                incluir m\'as regresores?
          \item ¿El aumento en $R^2$ justifica la complejidad adicional del
                modelo? (Use el criterio del $R^2$ ajustado).
        \end{enumerate}

  \item \textbf{Conclusiones y recomendaciones finales:} sintetice el
        an\'alisis completo.
        \begin{enumerate}[label=\alph*)]
          \item ¿Cu\'ales son los principales determinantes del puntaje de
                ingl\'es en Saber Pro para estudiantes de NI?
          \item ¿El modelo cumple razonablemente los supuestos de Gauss-Markov?
                ¿Qu\'e violaciones identific\'o?
          \item Si tuviera que recomendar un solo ajuste al modelo para
                mejorar su validez, ¿cu\'al ser\'ia? (Ejemplos: eliminar un
                regresor por VIF alto, usar errores robustos, excluir outliers,
                transformar variables).
          \item ¿El modelo tiene utilidad pr\'actica para predecir el
                rendimiento en Saber Pro de futuros estudiantes de NI?
        \end{enumerate}
\end{enumerate}

\vspace{10pt}
\rule{\linewidth}{0.4pt}
\begin{center}
{\small Cada entrega completa suma $+0.0625$ a la nota final.}
\end{center}

\end{document}
